\paragraph{Solvable and dynamical scaling-law models.}
Spectral and random-feature models derive power laws rather than only fitting them.  \citet{maloney2022solvable} connect power-law data statistics to parameter and sample scaling; \citet{bordelon2024dynamical} jointly analyze training time, width, and data; and \citet{zhang2024field} extends solvable random-feature scaling beyond the ridgeless limit.  The one-pass and data-reuse analyses of \citet{lin2024scaling,lin2025datareuse} isolate approximation, optimization, and effective-dimension terms.

\paragraph{Optimizer scaling and hyperparameter transfer.}
The $\mu$P framework stabilizes update statistics to enable hyperparameter transfer \cite{yang2021tensor}, and deep-linear theory studies transfer across width, depth, and data \cite{bordelon2025deep}.  Matrix-preconditioned optimizers also require scale-aware learning rates, blocking, normalization, and decay \cite{qiu2025hyperparameter}.  Contemporaneous studies fit optimizer-dependent laws in random-feature, quadratic, and representation-spectrum settings \cite{ramani2026optimizer,volkova2026robust,jha2026spectral}.  Our theory identifies visible second-moment spectrum as the basis-dependent mechanism.

\paragraph{Spectra of random-feature maps.}
Random-feature work studies statistical approximation, spectral concentration, and high-dimensional limits \cite{avron2017random,liao2018spectrum,ma2020slow,chen2022concentration,bosch2023precise}.  We instead study the diagonal of the feature covariance in optimizer coordinates.  Feature maps with similar eigenvalue spectra can expose different diagonals and hence different coordinatewise preconditioners; spectral-oracle controls separate information present in the representation from information accessible to a diagonal optimizer.
